% ===================================================================
%  GAMBAR No. 4 — Diwasa lan Tuwa.
%
%  Traduction, faite en 2026, en javanais d'aujourd'hui, sur l'ido de
%  texto/io. Regles de la colonne : voir 10-gambar-01.tex. Deux
%  scenes, deux numerotations : t04-c1-* la noce, t04-c2-*
%  l'anniversaire du grand-pere.
%
%  C'EST ICI QUE LA REGLE DE NIVEAU MORD POUR LA PREMIERE FOIS, ET
%  ELLE MORD EXACTEMENT COMME PREVU. Les tableaux 1 a 3 n'avaient que
%  des enfants et des maitres d'ecole ; celui-ci met en scene la
%  grand-mere (23) et le grand-pere (33), et la colonne passe au
%  NGOKO ALUS pour tout ce qui les concerne : gerah et non lara pour
%  la maladie, kersa et non arep pour la volonte, maos et non maca
%  pour la lecture, jumeneng et non ngadeg pour se lever, pirsa et
%  non weruh, ampeyan et non sikil pour la jambe, panjenengane et non
%  dheweke, kesupen et non lali. Le titre meme de la seconde scene le
%  porte : « Sowan », qui n'est pas « visiter » mais « se presenter
%  devant quelqu'un qu'on respecte », et qui ne se dit que dans ce
%  sens-la.
%
%  ET CELA VALIDE LA LISTE DE kolonoj.py, QUI N'ETAIT JUSQU'ICI
%  QU'UNE HYPOTHESE. Le ngoko alus ne change QUE les mots qui
%  denotent l'action, le corps ou le bien de la personne honoree ; la
%  grammaire, elle, reste en ngoko. « wis », « sing », « iki »,
%  « karo », « ora » ne bougent pas d'une lettre, meme dans la phrase
%  la plus deferente du tableau. C'est pourquoi la liste relevee —
%  sampun, ingkang, menika, kaliyan, mboten, kula — est la bonne :
%  ce sont tous des MOTS-OUTILS, et leur apparition ne signale pas
%  une politesse, elle signale que le niveau entier a glisse. Un
%  controle qui aurait releve dhahar ou gerah aurait crie a chaque
%  alinea de ce tableau-ci.
%
%  TROISIEME COUPLE NEERLANDAIS CONTRE PORTUGAIS, ET C'EST LE PLUS
%  NET DES TROIS. La fourchette (12) se dit porok en javanais, du
%  neerlandais vork ; elle se dit garpu en indonesien, du portugais
%  garfo. Apres potlot/pensil (tableau 1) et pit/sepeda (tableau 2),
%  la regle se voit : quand les deux langues ont emprunte le meme
%  objet, ce n'est presque jamais au meme preteur ni au meme siecle.
%  Le javanais prend au voisin qui est venu s'installer a Java ;
%  l'indonesien prend a celui qui passait sur la cote, puis a
%  l'ecole de la republique.
%
%  ET LA VERANDA (8) EN EST LA PREUVE INVERSE. L'en-tete du tableau 4
%  indonesien note que « beranda » lui vient du portugais varanda —
%  c'est meme cette ligne-la qui avait fini par reveler la couche
%  portugaise a cette colonne. Le javanais, lui, n'a rien emprunte du
%  tout : il dit emper, qui est son propre mot pour l'avancee de toit
%  sous laquelle on s'assied. Meme tableau, meme numero, meme chose,
%  et l'une des deux langues n'a eu besoin de personne.
%
%  LES DEGRES DE PARENTE SONT LE SECOND LEXIQUE QUE LE JAVANAIS NE
%  DOIT A PERSONNE, APRES LE CORPS (tableau 2). mbah kakung et mbah
%  putri les grands-parents, putu les petits-enfants, buyut
%  l'arriere-grand-pere, maratuwa les beaux-parents, mantu le gendre,
%  ipe le beau-frere, mbakyu la soeur ainee, pacangan les fiances,
%  randha la veuve, lola l'orphelin. Une famille se nomme dans la
%  langue de la maison, et la maison ne se traduit pas.
%
%  ET LE CIEL AUSSI EST JAVANAIS : rembulan la lune (17), lintang les
%  etoiles (18), peteng l'obscurite, geni le feu (30), wengi la nuit.
%  L'indonesien dit bulan, bintang, gelap, api, malam — pas un mot
%  commun. Comme au tableau 2, le lexique le plus vieux est celui qui
%  se ressemble le moins.
% ===================================================================

%%K t04-apar-1 apar
\VUsaut{4.0mm}
\VUtitre{15.0pt}{60}{GAMBAR No. 4}
\VUsaut{1.6mm}
\VUtitre{11.4pt}{0}{Diwasa lan Tuwa.}
\VUsaut{3.2mm}

%%K t04-tit-1 sub
\VUcentre{11.4pt}{0}{\textit{Adegan Kapisan.}}
\VUsaut{1.6mm}
\VUtitre{11.4pt}{0}{Pesta Mantenan.}
\VUsaut{2.6mm}

%%K t04-tit-2 sub
\VUpk{10.2pt}{40}{Mangsa gugur.}
\VUsaut{3.0mm}

%%K t04-c1-01-1 p
1. --- Para nom-noman wis padha gedhe. Salah sijine, Ioannes, omah-omah.
Kita melu \VUgras{pesta mantenan} kang nglumpukake kulawargane penganten
sakloron ing meja kang padha.

%%K t04-c1-02-1 p
2. --- Kita ana ing \VUgras{mangsa gugur,} ing \VUgras{sasi September.}
Angin adhem \VUgras{Oktober} lan \VUgras{November} durung ngrampasi
pedesan saka godhong ijone. Hawa sore isih anget lan wangi ; mula
bujana bengi iku dianakake ing sangisore \VUgras{emper}
\textsuperscript{(8)} kang madhep kebon.

%%K t04-c1-03-1 p
3. --- Ing kadohan, ing \VUgras{punthuk} \textsuperscript{(3)}, ana omahe
wong tuwane Ioannes, yaiku \VUgras{kastil} \textsuperscript{(1)} kang apik
lan kasingidan ing ijo-ijoan ; kebon anggur kang becik, kang ngasilake
anggur kang enak banget, nutupi ereng-erenge punthuk cilik iku. Tukang
kebon anggur kang ngolah lan ngopeni.

%%K t04-c1-04-1 p
4. --- Ing sadhuwure kebon anggur liwat sagerombolan \VUgras{ayam alas}
\textsuperscript{(2)} kang kaget dening tekane \VUgras{juru jaga alas}
\textsuperscript{(5)}. Wong iku, karo \VUgras{bedhil} ing keleke lan
\VUgras{tas buron} nyampir, njelajahi \VUgras{grumbulan}
\textsuperscript{(4)} karo \VUgras{asune} \textsuperscript{(6)}. Dheweke
njaga tanah iku lan ngalang-alangi wong mbebedhag nyolong supaya ora
ngentekake kewan buron.

%%K t04-c1-05-1 p
5. --- Ing sanjabane \VUgras{emper,} ana \VUgras{rambatan anggur} kang
menek, lan saka kono nggantung \VUgras{dhompolan anggur}
\textsuperscript{(7)} kang kandel.

%%K t04-c1-06-1 p
6. --- Kembang mangsa gugur kang endah, kang ngrengga \VUgras{meja}
\textsuperscript{(16)}, yaiku \VUgras{kembang krisan} \textsuperscript{(9)} ;
\VUgras{bapakne} \textsuperscript{(10)} penganten wadon nyelehake salah
sijine ing bolongan kancing jas frake.

%%K t04-c1-07-1 p
7. --- Bujana meh rampung. \VUgras{Pelayan} \textsuperscript{(11)} wiwit
nggawa piring-piring panganan panutup, lan sedhela maneh bakal
ngangkati piranti mangan kang wis ora dibutuhake maneh :
\VUgras{sendhok} \textsuperscript{(14)}, \VUgras{porok}
\textsuperscript{(12)}, \VUgras{ladhing} \textsuperscript{(13)} lan
\VUgras{piring gedhe} \textsuperscript{(15)}.

%%K t04-tit-3 sub
\VUpk{10.2pt}{40}{Sanak Kadang.}
\VUsaut{3.0mm}

%%K t04-c1-08-1 p
8. --- Kabeh katon bungah, lan esem \VUgras{ibune} \textsuperscript{(17)}
penganten lanang nalika nyawang anak-anake mbuktekake kabegjane.

%%K t04-c1-09-1 p
9. --- Mantenan iku nyawijekake kulawarga sugih loro ing tlatah kono.
Ioannes, \VUgras{penganten lanang} \textsuperscript{(18)}, yaiku
\VUgras{anake} juragan tanah kang gedhe, lan dheweke dadi
\VUgras{mantune} pengusaha kang trampil.

%%K t04-c1-10-1 p
10. --- \VUgras{Penganten} sakloron isih enom, nanging wis suwe padha
\VUgras{pacangan.} \VUgras{Wong wadon enom} \textsuperscript{(19)} iku,
Martha, ayu temen nganggo \VUgras{gaun putih} kang direngga kembang
jeruk.

%%K t04-c1-11-1 p
11. --- \VUgras{Maratuwa lanange} \textsuperscript{(20)} bungah banget
duwe \VUgras{mantu wadon} kang grapyak temen, lan \VUgras{maratuwa
wadone} \textsuperscript{(21)} Ioannes mesem nyawang \VUgras{anak
wadone} kang ayu temen.

%%K t04-c1-12-1 p
12. --- Ing pucuking meja, lungguh \VUgras{para tamu}
\textsuperscript{(22)} liyane : \VUgras{sanak kadang, kanca-kanca,
sedulur lanang lan sedulur wadon.} \VUgras{Kepala pelayan}
\textsuperscript{(23)}, karo serbet ing kelek, ngladeni wong-wong iku
kanthi cukat.

%%K t04-tit-4 sub
\VUpk{10.2pt}{40}{Kanca-kanca.}
\VUsaut{3.0mm}

%%K t04-c1-13-1 p
13. --- Ing sisih tengene meja, iki \VUgras{prawan} \textsuperscript{(24)},
\VUgras{kancane} penganten wadon enom iku, banjur sedulur lanange, kang
dadi \VUgras{pengiring penganten lanang} \textsuperscript{(25)}. Dheweke
ngobrol karo \VUgras{pengiring penganten wadone} \textsuperscript{(26)},
yaiku adhine wadon penganten lanang, kang lantaran mantenan iku dadi
\VUgras{ipe wadone} Martha. Bocah wadon iku ayu temen. Dheweke duwe
\VUgras{perhiasan} kang endah, \VUgras{kalung mutiara} lan
\VUgras{gelang} mas.

%%K t04-c1-14-1 p
14. --- Ing cedhake wong-wong iku, panjenengan kang ngadeg lan ngangkat
gelase yaiku \VUgras{kancane} \textsuperscript{(27)} kulawarga. Dheweke
salah sijine \VUgras{seksine penganten lanang.} Dheweke
\VUgras{ngucapake slamet,} ngombe kanggo \VUgras{kasarasan} penganten
enom sakloron lan nyuwunake kabegjan kang dawa. Dheweke nganggo
sandhangan upacara, jas frak karo \VUgras{sepatu lak}
\textsuperscript{(28)}. Dheweke dadi kanca lakune \VUgras{mbah putri}
\textsuperscript{(29)} kang wis \VUgras{randha.}

%%K t04-c1-15-1 p
15. --- Nom-noman kang nganggo \VUgras{jas frak} \textsuperscript{(31)},
kang brengose ireng tipis, yaiku \VUgras{ipe lanange}
\textsuperscript{(30)} Ioannes. Dheweke insinyur kang kajen. Dheweke
lungguh sandhinge \VUgras{nyonyah enom} \textsuperscript{(33)} kang apik
temen, yaiku \VUgras{mbakyune} Martha lan ibune bocah wadon cilik kang
teka nyekel tangane sedulur lanange, kanggo mènèhi kembang lan
\VUgras{muni} \VUgras{sugeng sonten} marang dheweke. Nyonyah iku duwe
\VUgras{anting} kang endah banget lan \VUgras{tutup sirah} kang apik.

%%K t04-c1-16-1 p
16. --- Sedulur lanang cilik iku, kang katon aneh temen amarga
\VUgras{bluse} \textsuperscript{(36)} lan \VUgras{kathok cendhake}
\textsuperscript{(37)} kang mekrok, mesakake banget. Dheweke
\VUgras{bocah lola} \textsuperscript{(35)}. Isih cilik banget dheweke wis
kelangan bapak lan ibune. Pakdhene dadi \VUgras{waline}
\textsuperscript{(38)}, lan ora omah-omah supaya bisa nggulawenthah
bocah kang mesakake iku. Dheweke wong kang loma. Rada
\VUgras{gundhul} lan \VUgras{lamur} banget ; dheweke nganggo
\VUgras{kacamata cepit.}

%%K t04-tit-5 sub
\VUsaut{2.4mm}
\VUpk{10.2pt}{40}{Panganan Panutup.}
\VUsaut{3.0mm}

%%K t04-c1-17-1 p
17. --- Ing mburine wong-wong kang bujana, ing meja kang katutupan
\VUgras{taplak,} \VUgras{panganan panutup} \textsuperscript{(39)} wis
disiyapake. Ing mangkok gedhe, iki woh-wohan : \VUgras{woh persik}
\textsuperscript{(40)}, \VUgras{woh pir} \textsuperscript{(41)},
\VUgras{apel} \textsuperscript{(42)}, \VUgras{aprikot}
\textsuperscript{(43)} lan \VUgras{anggur} \textsuperscript{(44)}. Ing
cedhake, \VUgras{nanas} \textsuperscript{(45)}, \VUgras{roti tar}
\textsuperscript{(46)} kang endah, \VUgras{botol minuman keras}
\textsuperscript{(48)} lan \VUgras{sampanye} \textsuperscript{(49)},
sarta tumpukan \VUgras{piring} \textsuperscript{(47)} kang resik.

%%K t04-c1-18-1 p
18. --- \VUgras{Juru gudhang anggur} \textsuperscript{(50)} mbukak
\VUgras{botol} \textsuperscript{(52)} anggur lawas nganggo
\VUgras{pembukak botol} \textsuperscript{(51)}. \VUgras{Sumbat gabuse}
\textsuperscript{(53)} katon angel temen dijupuk. Juru gudhang iku duwe
\VUgras{serbet} \textsuperscript{(54)} ing keleke.

%%K t04-c1-19-1 p
19. --- Sawise bujana bengi, para priya bakal menyang kamar udud, dene
para wanita bakal tata-tata kanggo pesta dansa.

%%K t04-tit-6 sub
\VUsaut{5.0mm}
\VUcentre{11.4pt}{0}{\textit{Adegan Kapindho.}}
\VUsaut{1.6mm}
\VUpk{11.4pt}{40}{Ulang taune Mbah Kakung.}
\VUsaut{2.6mm}

%%K t04-tit-7 sub
\VUpk{10.2pt}{40}{Sowan.}
\VUsaut{3.0mm}

%%K t04-c2-01-1 p
1. --- Sepuluh taun wis kliwat. Wong tuwane Ioannes wis sepuh lan
tresna banget marang putune telu, yaiku \VUgras{putu lanang}
\textsuperscript{(2)} loro lan \VUgras{putu wadon} \textsuperscript{(3)}
siji. Amarga dina iki \VUgras{ulang taune Mbah Kakung,} anak-anake teka
ngaturake sugeng ambal warsa.

%%K t04-c2-02-1 p
2. --- Petrus cilik isih enom banget ; umure lagi telung taun. Dheweke
weruh \VUgras{kucing} \textsuperscript{(1)} nyedhak. Dheweke kepengin
ngelus \VUgras{wulune} kang alus kaya sutra lan \VUgras{buntute} kang
dawa ; dheweke ora mikir yen ing sangisore \VUgras{sikile} kang luwes
ana \VUgras{kuku} landhep kang ndhelik, kang bisa uga nyakar dheweke.

%%K t04-c2-03-1 p
3. --- Mbakyune, Gabriela, kang nyekeli tangane, luwih tenanan, lan
Marcellus karo \VUgras{mantele} \textsuperscript{(4)} kang gedhe sarta
\VUgras{sabuk kulite} \textsuperscript{(5)} katon rada lanang, sanajan
kathok cendhake ngetokake \VUgras{kaos kakine} \textsuperscript{(6)}.

%%K t04-c2-04-1 p
4. --- Bapakne nganggo \VUgras{jas jabane} \textsuperscript{(7)} mangsa
adhem kang kandel karo \VUgras{krah wulune} \textsuperscript{(8)}, lan
ing \VUgras{kanthonge} \textsuperscript{(9)} dheweke nggawa kacu.
Bojone nganggo topi wol kang direngga \VUgras{wulu manuk}
\textsuperscript{(10)}, \VUgras{jakete} \textsuperscript{(11)},
\VUgras{sarung guluné wulu} \textsuperscript{(12)} lan \VUgras{sarung
tangan wulune} \textsuperscript{(13)}.

%%K t04-c2-05-1 p
5. --- Adhem banget, amarga mangsa adhem lagi ngrembaka.
\VUgras{Salju} \textsuperscript{(19)} tiba sedina muput lan payone
omah-omah putih kabeh. Bareng karo tekane \VUgras{wengi,} hawa adhem
saya nyocog. Ing langit, \VUgras{rembulan} \textsuperscript{(17)} lan
\VUgras{lintang} \textsuperscript{(18)} sumunar lan nembus
\VUgras{peteng.}

%%K t04-tit-8 sub
\VUsaut{4.2mm}
\VUpk{10.2pt}{40}{Lelara.}
\VUsaut{3.0mm}

%%K t04-c2-06-1 p
6. --- \VUgras{Wong wuta} \textsuperscript{(20)} kang mesakake, kang
liwat alun-alun ngarepe \VUgras{apotek} \textsuperscript{(16)} karo
dituntun \VUgras{asu pudele} \textsuperscript{(21)}, cilaka temen.
Iustinia, \VUgras{batur wadon} \textsuperscript{(14)}, nggawa saka pawon
\VUgras{mangkok} \textsuperscript{(15)} isi \VUgras{wedang} kang panas
banget lan diwenehi gula akeh.

%%K t04-c2-07-1 p
7. --- \VUgras{Mbah putri} \textsuperscript{(23)}, kang kersa nggoleki
\VUgras{kacamata gagange} \textsuperscript{(26)} ing meja supaya bisa
maos \VUgras{resep} \textsuperscript{(27)} kang lagi wae ditulis
\VUgras{dhokter} \textsuperscript{(25)}, jumeneng saka kursi malesé karo
nyanggong ing \VUgras{tekene} \textsuperscript{(24)}.

%%K t04-c2-08-1 p
8. --- Panjenengane kerep gerah \VUgras{lelara cilik, mumet, pilek,
lara untu} ; ampeyane wis ringkih lan mripate wis ora becik maneh.
Sanajan mangkono, panjenengane seneng nggeguyu \VUgras{dhoktere} kang
tuwa, kang tansah kepengin menehi \VUgras{obat cair}
\textsuperscript{(28)}. Dina iki panjenengane bungah banget amarga
kulawarga enome ana ing sakiwa-tengene.

%%K t04-c2-09-1 p
9. --- Kepenak temen ing kamar kang tinutup rapet. Ing \VUgras{tungku
marmer} \textsuperscript{(29)} murub \VUgras{geni kang endah}
\textsuperscript{(30)}, lan wong rumangsa begja ing \VUgras{padhang}
kang alus saka \VUgras{lampu} \textsuperscript{(31)} kang lagi wae
diurubake.

%%K t04-tit-9 sub
\VUsaut{4.2mm}
\VUpk{10.2pt}{40}{Mbah Kakung.}
\VUsaut{3.0mm}

%%K t04-c2-10-1 p
10. --- Malah \VUgras{Mbah Kakung} \textsuperscript{(33)} piyambak
kesupen yen panjenengane gerah, yen \VUgras{rematike} ngiket
panjenengane ing \VUgras{kursi malesé} \textsuperscript{(34)}, lan yen
wingi panjenengane isih \VUgras{gerah panas lan semaput.}
Panjenengane wis ora kengetan yen ing mangsa adhem kepungkur
panjenengane gerah \VUgras{radhang paru} lan \VUgras{watuk} serak kang
ngrekasakake.

%%K t04-c2-10-2 p
Nanging angel temen anggone waras. Saiki panjenengane rada mendhing.
Mung wae ampeyane, kang disanggongake ing \VUgras{bantal}
\textsuperscript{(38)} ing ndhuwur \VUgras{dhingklik}
\textsuperscript{(39)} cilik, isih krasa lara.

%%K t04-c2-11-1 p
11. --- Yen panjenengane wis kabungkus rapet ing \VUgras{klambi
kamare} \textsuperscript{(35)} kang anget karo \VUgras{kancing}
\textsuperscript{(36)} mutiara kang kandel, lan yen ing cedhake ana
\VUgras{bocah lola wadon} \textsuperscript{(40)} kang nganggo
\VUgras{kupluk} putih, \VUgras{gaun} \textsuperscript{(42)} biru lan
\VUgras{mantol cendhak} \textsuperscript{(41)}, kang macakake koran
kanggo panjenengane, panjenengane ora sambat.

%%K t04-c2-12-1 p
12. --- Saben taun, yen \VUgras{kalender} \textsuperscript{(43)} maneh
nuduhake \VUgras{tanggal} sepisan \VUgras{Desember,} panjenengane
ngentosi \VUgras{putu-putune} kanthi ora sabar, amarga panjenengane
pirsa yen putu-putune ora bakal kesupen marang \VUgras{tanggal} kang
wigati iku.

%%K t04-c2-13-1 p
13. --- Panjenengane kengetan kanthi trenyuh marang wektu biyen banget
nalika panjenengane piyambak sowan ngaturake sugeng marang
\VUgras{marsekal} kraton kang kaloka, kang \VUgras{potrete}
\textsuperscript{(44)} nggantung ing tembok, lan kang dadi
\VUgras{buyute} anak-anake.

\VUsaut{6.0mm}
\VUfilet{22mm}
